\problemname{Legacy Code}

Once again you lost days refactoring code, which never runs in the first place. Enough is enough -- your time is better spent writing a tool that finds unused code!

Your software is divided into packages and executables. A package is a collection of methods. Executables are packages defining among other methods exactly one method with name {\tt PROGRAM}. This method is executed on the start of the corresponding executable. Ordinary packages have no method named {\tt PROGRAM}.

Each method is uniquely identified by the combination of package and method names. E.g. the method with the identifier {\tt SuperGame::PROGRAM} would be the main method of the executable {\tt SuperGame}. 

For every method in your software you are given a list of methods directly invoking it. Thus you can easily identify methods, that are never called from any method. However, your task is more challenging: you have to find unused methods. These are methods that are never reached by the control flow  of any executable in your software.

\section*{Input}
The first line of the input contains an integer $N$, the number of methods in your software ($1\leq N\leq 400$). 

Each method is described by two lines, totaling in $2\cdot N$ lines. The first line consists of the unique identifier of the method and $k_i$, the number of methods directly invoking this one ($0 \leq k_i \leq N$). The second line consists of a set of $k_i$ identifiers of these calling methods or is empty if there are no such methods, i.e.~$k_i = 0$.

Method identifiers consist of a package name followed by two colons and a method name like {\tt Packagename::Methodname}. Both strings, the package and the method name, each consist of up to $20$ lowercase, uppercase characters or digits ({\tt a}-{\tt z}, {\tt A}-{\tt Z}, {\tt 0}-{\tt 9}).

There will be exactly $N$ different method identifiers mentioned in the input.

\section*{Output}
A line containing the number of unused methods in your software.
